\storyheader{Câu chuyện về cuộc sống bình dị \\ của chúng ta}{Ngày 04}{

Mắt tôi nheo nheo lại vì tiếp xúc quá lâu với màn hình máy tính, có vẻ
như chúng nó đang muốn được tôi thương tình mà thả lỏng để cho nhắm chặt
lại sau vài ngày thức trắng, cái lưng của tôi thì cũng trong tình thế
thảm hại không kém khi những cơn đau cứ liên tục nhói lên hành hạ cái
tâm trí tàn tạ của tôi.

Tuy nhiên mặc cho những điều đó thì hai tay của tôi lại đang gõ với một
tốc độ rất kinh khủng trên bàn phím, đến mức những con chữ trên màn hình
nhanh chóng lấp đầy từng dòng một. Tôi mặc xác tình trạng đáng báo động
hiện tại của bản thân và tiếp tục viết, hiện tại là 23:45, hạn chót của
tôi là mười hai giờ nên về cơ bản thì hiện tại tôi đang ở trong tình thế
hết sức nguy hiểm.

Từng phút, từng phút một trôi qua, tôi cảm thấy như thể sự nghiệp và ước
mơ như đang dần tan biến khỏi bàn tay và hóa thành cát bụi bay đi theo
cơn gió, bỏ tôi ở lại với những tháng ngày thất nghiệp và ăn bám vợ. Và
tất nhiên tôi hoàn toàn không thể để cái tình huống đó xảy ra được rồi,
nên là tôi vẫn liên tục ép bản thân của mình vượt quá giới hạn.

Cố lên nào cái lưng, cánh tay và đôi mắt của tao ơi! Mày rất mạnh! Mày
đã từng trong tình thế còn căng thẳng hơn hồi đại học khi làm bài thuyết
trình cơ mà! Hãy nhớ lại những giọt mồ hôi, từng sự buồn ngủ, từng cơn
đau cơ và cách bọn mày đã vượt qua đi! Hãy cho đám nhà sản xuất thấy sức
mạnh của chúng ta đi nào!

Được rồi, ngưng ảo tưởng cái trò sức mạnh tình bạn và tiếp tục làm việc
thôi, sau khoảng vài phút gõ chữ như một cái máy thì cuối cùng, cuối
cùng tôi cũng đã hoàn thành chương mới của mình. Nhẹ nhõm nằm phịch
xuống, hai con mắt của tôi lúc này bắt đầu biểu tình cho sự nghỉ ngơi
của mình, thế nhưng tôi chưa thể ngủ được, phải tận hưởng cái dư vị
chiến thắng này mới được!

Hỡi những kẻ trong nhà sản xuất, hãy nhìn đi! Hãy nhìn vào cái kẻ mà các
người cho rằng sẽ trễ hạn deadline đi! Nhìn thấy chưa hả? Đây chính là
sức mạnh của niềm đam mê và khát vọng đấy! Hahahahahahaha!

Trong lúc đang tận hưởng cái cảm giác chiến thắng của mình thì đột nhiên
tôi cảm thấy một lon nước lạnh áp vào má mình.

``Chị vất vả rồi, chị Yuzuha.'' Cô gái ấy mỉm cười với tôi.

Tôi nhận lấy lon nước từ em ấy. Mỗi ngày chỉ cần em ấy nở một nụ cười
với tôi thôi thì cảm tưởng như mọi mệt mỏi đều sẽ tan biến hết, áp lực
do đám nhà sản xuất ấy tạo ra làm sao mà so được với động lực mà mỗi
ngày em ấy tiếp thêm cho tôi cơ chứ!

Quên giới thiệu, đây chính là Mirai, hay nói thẳng ra là người vợ yêu
dấu của tôi. Tôi là tiền bối của em ấy ở trong câu lạc bộ văn học thời
còn đi học. Em ấy đúng là kiểu một người vợ vô cùng đảm đang và tài giỏi
mà bất kỳ ai cũng mong muốn có được khi mà em ấy quán xuyến gần như mọi
công việc trong nhà thay cho một đứa tiểu thuyết gia hoàn toàn không có
khả năng sống tự lập như tôi. Em ấy còn là giáo viên dạy văn nữa.

Để tránh bị hiểu lầm, tuy tôi thường gây khá nhiều phiền phức, không hay
giúp đỡ em ấy trong việc nhà và chỉ ngồi lì trong phòng nhưng tôi vẫn là
một tiểu thuyết gia nổi tiếng và thu nhập của tôi rất cao, thế nên cũng
không phải tôi ăn bám em ấy hoàn toàn đâu.

Tuy hơi khó tin, nhưng Mirai là người đã tỏ tình tôi trước vào năm ba,
khi đó thì tôi đã sắp tốt nghiệp rồi và chuẩn bị giao lại câu lạc bộ văn
học cho Mirai thì em ấy đã kéo tôi ra sau trường. Lúc đó em ấy khóc nức
nở khiến tôi thực sự rất là hoảng luôn đó! Rồi sau đó em ấy nhào đến ôm
tôi, vừa khóc vừa nói. ``Em yêu chị, hội trưởng! Em đã yêu chị từ rất
lâu rồi!''

Đại loại thế đấy, sau đó thì tôi với em ấy ra mắt gia đình hai bên, rồi
sau khi tốt nghiệp đại học thì Mirai nhanh chóng trở thành giáo viên,
tôi thì dự thi một cái giải viết truyện ngắn và được giải nhất, nhờ sự
động viên của Mirai mà tôi bắt đầu viết truyện dài và gửi cho một nhà
sản xuất. Bất ngờ là nó lại được duyệt và được đón nhận nồng nhiệt, cuối
cùng thì có tôi, một tiểu thuyết gia đình đám.

``Chà, cuối cùng chị cũng hoàn thành nhỉ?'' Mirai ngồi xuống kế tôi và
bắt chuyện. ``Em còn nhớ rõ chị dám cả gan ăn mì ly và hốc toàn nước
tăng lực trong một ngày đấy.''

``Chuyện bất khả kháng mà em!'' Tôi khui lon nước ra và chống chế cho
bản thân. ``Chị cũng muốn ăn cơm với em lắm, nhưng deadline không cho
phép!''

``Lại còn từ chối tắm chung nữa!'' Mirai làm vẻ hờn dỗi và quay mặt đi.
``Lúc đó em thật sự đã cô đơn lắm đó!''

Thôi chết tôi! Vì cái chương mà tôi bị vướng là chương cuối của một cái
truyện được độc giả rất yêu thích nên bên nhà sản xuất làm khá căng,
thành thử ra là tôi khóa chặt cửa phòng làm việc lại để tránh bị mất tập
trung, trữ hàng đống mì ly lẫn nước tăng lực để mà sinh tồn trong này.

Và, trong mấy ngày đó thì cô vợ yêu quý của tôi đã phải ăn cơm một mình,
tắm một mình, và cả ngủ một mình nữa. Chết tiệt cái đám nhà sản xuất đó!
Chương cuối bị dời thì có làm sao đâu chứ! Dăm ba mấy tên độc giả mà lại
đòi quan trọng hơn vợ của tôi sao?

Nhìn em ấy vẫn đang phồng má đầy giận dỗi, mắt cứ liếc liếc tôi như thể
đang mong chờ thứ gì đó. Mong chờ thứ gì đó ấy à? Nếu như để dỗ dành một
người đang cô đơn và thiếu hơi người yêu, thì chỉ cần bù đắp bằng cái
hơi đấy thôi nhỉ?

Tôi nắm lấy hai vai của Mirai, đẩy em ấy nằm xuống và nằm đè lên.

``Ế?'' Mirai ngước lên nhìn tôi với một khuôn mặt ngơ ngác trông như
không hiểu tôi đang định làm gì.

Tuy nhiên, tôi đã kết hôn với Mirai quá lâu để có thể đọc ra hoàn toàn
cảm xúc bên trong ánh mắt của em ấy, và hiện tại có vẻ như em ấy muốn
tôi tiếp tục nhỉ? Nghĩ vậy, tôi luồn tay vào áo của Mirai, đưa miệng sát
lại gần và thì thầm vào tai của em ấy.

``Xin lỗi vì đã khiến em phải cô đơn nhé. Chị sẽ bù đắp cho em ngay bây
giờ.''

``Ưm!'' Mirai bật ra tiếng rên khe khẽ.

Cơ thể của Mirai rất mảnh mai và nhỏ nhắn luôn khiến tôi có một cảm giác
gì đó khiến mình buộc phải bảo vệ, nâng niu và nhẹ nhàng với em ấy.
Nhưng với tình thế hiện tại, khi mà em ấy đang bày ra một khuôn mặt đỏ
bừng, ánh mắt rơm rớm nước vì ngại ngùng nhưng hai cánh tay vẫn nằm yên
như thể ngầm bảo tôi hãy tiếp tục việc này.

Tôi muốn nữa, tôi muốn tiếp tục nhìn thấy Mirai như thế này, tôi muốn
nhìn ngắm khuôn mặt ngại ngùng của em ấy, tôi muốn thấy phản ứng của cơ
thể mảnh mai này, tôi muốn nghe tiếng rên nhỏ nhẹ và ngọt ngào ấy.

Aaaa, gian lận, em gian lận quá đấy Mirai, nếu cứ đáng yêu như vậy thì
làm sao chị có thể ngừng việc này lại đây?

``Đ-được rồi! Em không giận nữa đâu!'' Mirai đột nhiên đẩy tôi ra. ``Chị
mau đi tắm đi!''

``Ơ?''

``Và!'' Mirai chỉ thẳng vào mặt tôi. ``Hôm nay không có làm!''

*

``Để xem nên chết kiểu nào thì ít phiền đến Mirai nhất ấy nhỉ?'' Tôi lẩm
bẩm một mình trong phòng tắm.

Không không không, nếu mà chết thì làm sao tôi được tiếp tục tận hưởng
cuộc sống hôn nhân vô cùng hạnh phúc và viên mãn này đây chứ? Mặc dù bây
giờ tôi có lẽ là bị Mirai ghét rồi, nhưng mà tôi vẫn muốn tiếp tục ở bên
em ấy cơ!

Nhưng mà, quả nhiên là Mirai chán ghét tôi rồi! Cũng phải thôi, là do
tôi đã tệ bạc bỏ mặc em ấy một mình suốt tận mấy ngày liền cơ mà! Bây
giờ có khi em ấy ghét cái bản mặt của tôi lắm rồi ấy chứ! Vậy mà vẫn cất
công đem nước vào cho tôi, Mirai thật quá tốt bụng mà!

Tôi giãy một hồi lâu trong bồn tắm thì cũng bắt đầu thấy chóng mặt, nghĩ
rằng mình đã ngâm nước quá lâu nên tôi đi ra ngoài, mặc quần áo và vác
cả cái bộ mặt trầm cảm vì mấy cái suy nghĩ tiêu cực vừa nãy.

``Chị xong rồi đây...'' Thế nhưng trước mắt tôi, Mirai đang sửa soạn mền
gối trên ghế sofa.

``A, chị Yuzuha?''

Em ấy... lại định ngủ ở đó sao? Dù cho tôi đã xong hết công việc?

``...Tại sao em lại muốn ngủ riêng?'' Tôi tiến đến hỏi Mirai, và nắm lấy
hai vai của em ấy. ``Nếu chị làm em cô đơn, chị thành thật xin lỗi, vì
tính chất công việc nên là chị không thể dành toàn bộ thời gian bên em.
Nếu em muốn chia ta-''

``Từ từ! Chị bi quan quá rồi đấy! Em yêu chị còn không hết, sao mà lại
muốn chia tay được!'' Mirai vội vàng bịt cái mồm đang dần đưa ra một
viễn cảnh siêu tồi tệ của tôi. ``Em không có ghét chị, cũng không giận
chị vì chị bận rộn với công việc đâu!''

``Nếu vậy thì, tại sao em lại muốn ngủ riêng?'' Tôi nói với cái miệng
đang bị em ấy bịt lại. ``Lại còn nói không muốn làm chuyện đó nữa...''

``Bởi vì...'' Mirai ngập ngừng nói với tôi, khuôn mặt đỏ đến tận mang
tai vì xấu hổ. ``Với thời gian thiếu hình bóng và hơi ấm của chị, em sợ
khi ngủ chung thì sẽ không kiềm chế được, và nếu bắt đầu thì sẽ không
thể dừng lại mất...''

``Vậy nên em mới không muốn làm?''

``Chị vừa mới làm việc rất vất vả, em không muốn vì sự ích kỷ của mình
khiến chị mệt thêm.'' Mirai lấy hai tay che khuôn mặt đã đỏ bừng của
mình. ``Em xin lỗi nếu việc đó khiến chị buồn, em vẫn rất yêu chị
mà...''

Ra vậy, Mirai vì thấy tôi làm việc vất vả nên là mới không muốn làm
chuyện đó nhỉ? Và vì đã mấy ngày rồi không được tiếp xúc với tôi nên em
ấy không thể chắc chắn bản thân có thể kiềm chế nên đã quyết định ngủ
riêng để tôi được thư giãn. Vậy mà tôi lại nghĩ Mirai ghét mình, đúng là
ngốc thật.

``Ngốc à.'' Tôi quàng tay ôm lấy cơ thể mảnh mai và nhỏ nhắn ấy của
Mirai. ``Liều thuốc bổ dưỡng nhất của chị chính là em đó.''

Sau đó, tôi bế Mirai lên theo kiểu công chúa, cơ thể Mirai thật sự rất
nhẹ, và em ấy gần như nằm trọn trong lòng tôi.

``A, khoan đã chị Yuzuha!'' Mirai vùng vẫy trên tay tôi. ``Chị thật sự
cần được nghỉ ngơi mà!''

``Quên nói với em, vì chị đã viết xong truyện rồi nên là sẽ có được một
kỳ nghỉ dài hạn.'' Tôi dừng lại ở trước cửa phòng ngủ vì nhớ ra một
chuyện quan trọng, trên tay vẫn đang bế Mirai và vác đống mền gối của em
ấy.

``Thật sao?''

``Vậy, chúng ta ngủ cùng với nhau nhé?'' Tôi đặt lên đôi môi mềm mại của
em ấy một nụ hôn nhẹ.

``Đồ ngốc...'' Mirai đỏ bừng mặt quay sang chỗ khác. Tôi xem đó là một
lời đồng ý đáng yêu của em ấy nên cũng xuôi theo.

Lại một ngày nữa trôi qua.

Một ngày rất đỗi bình thường, trong câu chuyện về cuộc sống bình dị và
tràn đầy hạnh phúc của cả hai chúng tôi.
}